\begin{proof}[Proof of \cref{obj:lem:honest-interval-construction}]
Fix the parameters in the statement. For every observed sample \(o\in O^n\), write the endpoints from \cref{obj:def:mm-interval} as
\[
L(o)=\max\{-1,\widehat\tau_{\mathrm{MM}}(o)-h_{n,d,q,\alpha}\},
\qquad
U(o)=\min\{1,\widehat\tau_{\mathrm{MM}}(o)+h_{n,d,q,\alpha}\}.
\]
The radius satisfies
\[
0\le h_{n,d,q,\alpha}
=\min\left\{2,\sqrt{\frac{C_Ur_{n,d,q}}{\alpha}}\right\}
\le2.
\]

\begin{enumerate}
\item The endpoint formulas give
\[
L(o)\ge-1,\qquad U(o)\le1,
\]
and
\[
L(o)\ge\widehat\tau_{\mathrm{MM}}(o)-h_{n,d,q,\alpha},
\qquad
U(o)\le\widehat\tau_{\mathrm{MM}}(o)+h_{n,d,q,\alpha}.
\]
Subtracting yields both bounds
\[
U(o)-L(o)\le2,
\qquad
U(o)-L(o)\le2h_{n,d,q,\alpha}.
\]
Consequently,
\[
U(o)-L(o)\le\min\{2,2h_{n,d,q,\alpha}\}.
\]
If \(2\le\sqrt{C_Ur_{n,d,q}/\alpha}\), then the radius equals \(2\), and both sides of
\[
\min\{2,2h_{n,d,q,\alpha}\}
=
\min\left\{2,2\sqrt{\frac{C_Ur_{n,d,q}}{\alpha}}\right\}
\]
equal \(2\). If \(\sqrt{C_Ur_{n,d,q}/\alpha}\le2\), the identity follows by substituting that square root for the radius. This proves the pointwise length bound and its displayed equality.
% lean: CausalSmith.Stat.MarNearcompleteFrontier.honest_interval_construction

\item We establish the target range used in the coverage argument. For a fixed \(P\in\mathcal M(d,q)\), define
\[
\mathcal W_P=\operatorname{supp}(P),
\qquad
p_w=P(\{w\}),
\qquad
\Delta_w=Y(1)(w)-Y(0)(w)
\quad(w\in\mathcal W_P).
\]
The support is finite, and the binary potential outcomes give
\[
p_w\ge0,\qquad
\sum_{w\in\mathcal W_P}p_w=1,
\qquad
-1\le\Delta_w\le1.
\]
Thus
\[
-p_w\le p_w\Delta_w\le p_w.
\]
Summing these inequalities proves
\[
-1
=-\sum_{w\in\mathcal W_P}p_w
\le
\tau(P)=\sum_{w\in\mathcal W_P}p_w\Delta_w
\le
\sum_{w\in\mathcal W_P}p_w
=1.
\]
% lean: CausalSmith.Stat.MarNearcompleteFrontier.tau_range

\item The estimator also satisfies
\[
-1\le\widehat\tau_{\mathrm{MM}}(o)\le1
\qquad(o\in O^n).
\]
To derive this from \cref{obj:def:mm-estimator}, first observe that, for every real \(t\),
\[
-1\le\Pi_{[-1,1]}(t)=\max\{-1,\min\{1,t\}\}\le1:
\]
the lower inequality follows from the maximum, and the upper inequality follows because both arguments of that maximum are at most \(1\). This establishes the range when \(q=1\). When \(q\ne1\) and
\(\ell_n<128\) or \(d\ge n\ell_n\), the fallback branch is likewise a projection and has the same range.

For the remaining branch, define the Poisson weights and conditional stream averages by
\[
\omega_N=e^{-n/2}\frac{(n/2)^N}{N!}
\quad(N\in\mathbb N_0),
\]
\[
\overline\tau_N(o)
=
\frac{1}{4^N}
\sum_{\boldsymbol\zeta\in\{1,2,3,4\}^N}
\widetilde\tau(o;N,\boldsymbol\zeta)
\quad(0\le N\le n,\ o\in O^n),
\]
where \(\widetilde\tau(o;N,\boldsymbol\zeta)\) is the projected randomized estimate in \cref{obj:def:mm-estimator} with count \(N\) and stream assignment \(\boldsymbol\zeta\). There are \(4^N\) assignments, including one empty assignment when \(N=0\), and each estimate lies in \([-1,1]\). Therefore
\[
-4^N
\le
\sum_{\boldsymbol\zeta\in\{1,2,3,4\}^N}
\widetilde\tau(o;N,\boldsymbol\zeta)
\le4^N,
\qquad
-1\le\overline\tau_N(o)\le1.
\]
Moreover,
\[
\omega_N\ge0,\qquad
\sum_{N=0}^{\infty}\omega_N
=e^{-n/2}e^{n/2}=1,
\qquad
\sum_{N=0}^{n}\omega_N\le1.
\]
Since the randomized estimate equals zero when \(N>n\), this branch satisfies
\[
\widehat\tau_{\mathrm{MM}}(o)
=\sum_{N=0}^{n}\omega_N\overline\tau_N(o).
\]
Multiplying the stream-average bounds by the nonnegative weights and summing gives
\[
-1
\le-\sum_{N=0}^{n}\omega_N
\le\widehat\tau_{\mathrm{MM}}(o)
\le\sum_{N=0}^{n}\omega_N
\le1.
\]
% lean: CausalSmith.Stat.MarNearcompleteFrontier.tauhatMM_range

In particular, if \(h_{n,d,q,\alpha}=2\), the target and estimator ranges imply
\[
\widehat\tau_{\mathrm{MM}}(o)-2\le-1\le\tau(P),
\qquad
\tau(P)\le1\le\widehat\tau_{\mathrm{MM}}(o)+2.
\]
Taking the maximum in the lower endpoint and the minimum in the upper endpoint yields
\[
L(o)\le\tau(P)\le U(o)
\qquad(P\in\mathcal M(d,q),\ o\in O^n).
\]
% lean: CausalSmith.Stat.MarNearcompleteFrontier.honest_interval_construction

\item Suppose now that \(h_{n,d,q,\alpha}<2\), and fix \(P\in\mathcal M(d,q)\). Define the sample probability measure and squared-error function by
\[
\mu_{\mathrm{obs}}=P_O^{\otimes n},
\qquad
\mathcal E_P(o)
=\bigl(\widehat\tau_{\mathrm{MM}}(o)-\tau(P)\bigr)^2
\quad(o\in O^n).
\]
If \(\sqrt{C_Ur_{n,d,q}/\alpha}\le2\), the radius definition gives
\(h_{n,d,q,\alpha}=\sqrt{C_Ur_{n,d,q}/\alpha}\).
If that square root exceeds \(2\), the radius equals \(2\), contradicting the present case. Hence
\[
h_{n,d,q,\alpha}
=\sqrt{\frac{C_Ur_{n,d,q}}{\alpha}},
\qquad
h_{n,d,q,\alpha}^2=\frac{C_Ur_{n,d,q}}{\alpha}.
\]
By \cref{obj:def:rate} and \(n\ge1\),
\[
r_{n,d,q}\ge\frac1n>0.
\]
Together with \(C_U>0\) and \(\alpha>0\), this proves
\[
h_{n,d,q,\alpha}>0,
\qquad
h_{n,d,q,\alpha}^2>0,
\qquad
C_Ur_{n,d,q}=\alpha h_{n,d,q,\alpha}^2.
\]

Define the noncoverage and squared-error tail events by
\[
\mathcal B_P=\{o\in O^n:\tau(P)\notin[L(o),U(o)]\},
\qquad
\mathcal T_P=\{o\in O^n:h_{n,d,q,\alpha}^2\le\mathcal E_P(o)\}.
\]
We claim that
\[
\mathcal B_P\subseteq\mathcal T_P.
\]
Indeed, noncoverage gives either \(\tau(P)<L(o)\) or \(U(o)<\tau(P)\). In the first case, if
\(-1\le\widehat\tau_{\mathrm{MM}}(o)-h_{n,d,q,\alpha}\), then
\[
L(o)=\widehat\tau_{\mathrm{MM}}(o)-h_{n,d,q,\alpha},
\qquad
h_{n,d,q,\alpha}
<\widehat\tau_{\mathrm{MM}}(o)-\tau(P).
\]
Otherwise \(L(o)=-1\), contradicting \(\tau(P)\ge-1\). In the second case, if
\(\widehat\tau_{\mathrm{MM}}(o)+h_{n,d,q,\alpha}\le1\), then
\[
U(o)=\widehat\tau_{\mathrm{MM}}(o)+h_{n,d,q,\alpha},
\qquad
h_{n,d,q,\alpha}
<\tau(P)-\widehat\tau_{\mathrm{MM}}(o).
\]
Otherwise \(U(o)=1\), contradicting \(\tau(P)\le1\). Squaring the positive inequalities in the two possible cases gives
\[
h_{n,d,q,\alpha}^2
\le\bigl(\widehat\tau_{\mathrm{MM}}(o)-\tau(P)\bigr)^2,
\]
as claimed.

The sample space is finite, so the squared-error function is integrable. Its nonnegativity and Markov's inequality give
\[
h_{n,d,q,\alpha}^2\,\mu_{\mathrm{obs}}(\mathcal T_P)
\le
\int\mathcal E_P(o)\,d\mu_{\mathrm{obs}}(o).
\]
Applying the assumed uniform risk certificate to this iid sample law, and using event inclusion, yields
\[
h_{n,d,q,\alpha}^2\,\mu_{\mathrm{obs}}(\mathcal B_P)
\le
h_{n,d,q,\alpha}^2\,\mu_{\mathrm{obs}}(\mathcal T_P)
\le
\int\mathcal E_P(o)\,d\mu_{\mathrm{obs}}(o)
\le
C_Ur_{n,d,q}
=\alpha h_{n,d,q,\alpha}^2.
\]
Dividing by the strictly positive squared radius proves
\[
P_O^{\otimes n}\{\tau(P)\notin\widehat I_{\mathrm{MM},\alpha}(O^n)\}
\le\alpha.
\]
% lean: CausalSmith.Stat.MarNearcompleteFrontier.honest_interval_construction

\item To obtain honesty, fix any \(P\in\mathcal M(d,q)\) and split according to whether \(h_{n,d,q,\alpha}=2\). In that case, the pointwise inclusion proved above gives
\[
P_O^{\otimes n}\{\tau(P)\in\widehat I_{\mathrm{MM},\alpha}(O^n)\}
=1\ge1-\alpha.
\]
In the other case, \(h_{n,d,q,\alpha}\le2\) implies \(h_{n,d,q,\alpha}<2\). The coverage event is measurable, and the preceding noncoverage bound therefore gives
\[
P_O^{\otimes n}\{\tau(P)\in\widehat I_{\mathrm{MM},\alpha}(O^n)\}
=
1-P_O^{\otimes n}\{\tau(P)\notin\widehat I_{\mathrm{MM},\alpha}(O^n)\}
\ge1-\alpha.
\]
Since \(P\) was arbitrary, the interval has the asserted uniform coverage.
% lean: CausalSmith.Stat.MarNearcompleteFrontier.honest_interval_construction

\item Finally, the estimator range and nonnegative radius ensure that the endpoints form a nonempty interval, so
\[
0\le\operatorname{length}\{\widehat I_{\mathrm{MM},\alpha}(o)\}
=U(o)-L(o).
\]
These functions are integrable on the finite sample space. Integrating the pointwise bound from the first step under the probability measure \(P_O^{\otimes n}\) gives
\[
\begin{aligned}
\mathbb E_{P_O^{\otimes n}}
\!\left[\operatorname{length}\{\widehat I_{\mathrm{MM},\alpha}(O^n)\}\right]
&\le
\int\min\{2,2h_{n,d,q,\alpha}\}\,dP_O^{\otimes n}\\
&=\min\{2,2h_{n,d,q,\alpha}\}\\
&=\min\left\{2,2\sqrt{\frac{C_Ur_{n,d,q}}{\alpha}}\right\}.
\end{aligned}
\]
This holds for every \(P\in\mathcal M(d,q)\).
% lean: CausalSmith.Stat.MarNearcompleteFrontier.honest_interval_construction
\end{enumerate}
\end{proof}
